
\documentclass[12pt,a4paper]{article}

\usepackage[a4paper,margin=1in]{geometry}
\usepackage{graphicx}
\usepackage{booktabs}
\usepackage{amsmath,amssymb}
\usepackage{float}
\usepackage{hyperref}
\usepackage{xcolor}
\usepackage{listings}
\usepackage{enumitem}
\usepackage{fancyhdr}
\usepackage{longtable}

\hypersetup{colorlinks=true,linkcolor=blue,urlcolor=blue}

\pagestyle{fancy}
\fancyhf{}
\lhead{ICS1512}
\rhead{Machine Learning Laboratory}
\cfoot{\thepage}

\lstset{
language=Python,
basicstyle=\ttfamily\small,
numbers=left,
frame=single,
breaklines=true,
showstringspaces=false
}

\begin{document}

\begin{tabular}{|p{4cm}|p{11cm}|}
\hline
Experiment No. & 6 \\ \hline
Experiment Title & Dimensionality Reduction and Model Evaluation (With and Without PCA)
\footnote{\textbf{GitHub Repository: }\url{PASTE_YOUR_GITHUB_LINK_HERE}}\\ \hline
Student Name & Samyuktha V \\ \hline
Register Number & 3122247001053 \\ \hline
Faculty & Dr. Poreddy Ajay Kumar Reddy \\ \hline
Submission Date & 12/08/26 \\ \hline
\end{tabular}

\section{Objective}
To study the effect of dimensionality reduction using Principal Component Analysis
(PCA) on the performance of 10 machine learning classifiers -- training and
validating each model both without PCA (original feature space) and with PCA
(reduced feature space), with hyperparameter tuning and 5-fold cross-validation
applied in both settings.

\section{Problem Statement}
Given 30 numerical features describing breast mass characteristics, classify each
sample as Malignant or Benign, and determine whether reducing those 30 features to
a smaller PCA-derived feature set helps, hurts, or has no meaningful effect on
each of 10 different classifiers' accuracy, F1-score, and cross-validation
stability.

\section{Dataset Description}
\begin{longtable}{|p{6cm}|p{8cm}|}
\hline
Dataset Source & UCI Machine Learning Repository (Wisconsin Diagnostic Breast
Cancer Dataset, also available via \texttt{sklearn.datasets.load\_breast\_cancer()}) \\ \hline
Number of Samples & 569 \\ \hline
Number of Features & 30 numerical attributes \\ \hline
Target Classes \& Distribution & 2 classes -- Benign (357, 62.7\%), Malignant
(212, 37.3\%) \\ \hline
Missing Values & 0 -- no imputation required \\ \hline
Preprocessing & Standardization (\texttt{StandardScaler}, fit on train only);
no categorical encoding needed (all features numeric) \\ \hline
Train-Test Split & 80\% train (455 samples) / 20\% test (114 samples), stratified,
\texttt{random\_state=42} \\ \hline
\end{longtable}

\section{PCA Design Choice}

\begin{longtable}{|p{5cm}|p{4cm}|p{3cm}|p{3cm}|}
\hline
Setting & Chosen Components / Variance Target & Explained Variance (\%) & Justification\\ \hline
With-PCA & 10 components (target: 95\% cumulative variance) & 95.27\% & Smallest
number of components that reaches the chosen 95\% variance threshold, balancing
information retention against dimensionality reduction (30 $\rightarrow$ 10
features, a 67\% reduction) \\ \hline
\end{longtable}

\begin{figure}[H]
\centering
\includegraphics[width=0.7\linewidth]{scree_plot.png}
\caption{PCA Scree Plot -- cumulative explained variance vs. number of components}
\end{figure}

\textbf{Inference:} the first 10 principal components capture 95.27\% of the total
variance in the original 30 standardized features, meaning the remaining 20
components together contribute under 5\% -- most of the dataset's real
information content is concentrated in a small number of directions, largely due
to the strong correlations between related features (e.g. radius, perimeter, and
area measurements) observed during EDA.

\section{Machine Learning Algorithms}

Ten classifiers were evaluated: Support Vector Machine (RBF kernel), Gaussian
Naive Bayes, K-Nearest Neighbors, Logistic Regression, Decision Tree, Random
Forest, AdaBoost, Gradient Boosting, XGBoost, and a Stacked Ensemble (SVM +
Naive Bayes + Decision Tree base learners, Logistic Regression meta-learner).
Working principles and mathematical formulations for Logistic Regression, SVM,
Naive Bayes, KNN, and the ensemble methods (Bagging/Boosting/Stacking family)
are covered in the Experiment 2/4/7 reports; this experiment focuses on how
PCA-reduced input features affect each of them rather than re-deriving their
individual theory.

\section{Experimental Setup}
\begin{tabular}{ll}
\toprule
Python Version & 3.11 \\
Scikit-Learn Version & 1.9.0 \\
Additional Library & XGBoost 3.4.1 \\
IDE & Jupyter Notebook \\
Operating System & Windows \\
Hardware & Intel i5, 16GB RAM \\
\bottomrule
\end{tabular}

\section{Hyperparameter Grids Defined}

\begin{longtable}{ll}
\toprule
Model & Search Space\\
\midrule
SVM & kernel=rbf; C $\in \{1, 10\}$; $\gamma \in \{$scale, auto$\}$ \\
Naive Bayes & var\_smoothing $\in \{10^{-9}, 10^{-8}, 10^{-7}\}$ \\
KNN & $k \in \{3, 5, 7\}$; weights $\in \{$uniform, distance$\}$ \\
Logistic Regression & $C \in \{0.1, 1, 10\}$ \\
Decision Tree & max\_depth $\in \{3, 5,$ None$\}$; min\_samples\_split $\in \{2, 5\}$ \\
Random Forest & n\_estimators $\in \{50, 100\}$; max\_depth $\in \{5,$ None$\}$ \\
AdaBoost & n\_estimators $\in \{50, 100\}$; learning\_rate $\in \{0.1, 1.0\}$ \\
Gradient Boosting & n\_estimators $\in \{50, 100\}$; learning\_rate $\in \{0.1, 1.0\}$; max\_depth $\in \{2, 3\}$ \\
XGBoost & n\_estimators $\in \{50, 100\}$; learning\_rate $\in \{0.1, 1.0\}$; max\_depth $\in \{3, 5\}$ \\
Stacking & Fixed structure (base learners + meta-learner listed above), not tuned \\
\bottomrule
\end{longtable}

Each model was tuned with \texttt{GridSearchCV} (5-fold CV, scoring = accuracy)
separately on the No-PCA (30-feature) and With-PCA (10-feature) training data.

\section{Results}

\subsection*{Best Hyperparameters Found}
\begin{longtable}{lll}
\toprule
Model & Best Params (No-PCA) & Best Params (With-PCA)\\
\midrule
SVM & C=10, gamma=scale, kernel=rbf & C=1, gamma=scale, kernel=rbf \\
Naive Bayes & var\_smoothing=1e-9 & var\_smoothing=1e-9 \\
KNN & n\_neighbors=3, weights=uniform & n\_neighbors=7, weights=uniform \\
Logistic Regression & C=0.1 & C=0.1 \\
Decision Tree & max\_depth=5, min\_samples\_split=5 & max\_depth=5, min\_samples\_split=2 \\
Random Forest & max\_depth=None, n\_estimators=100 & max\_depth=None, n\_estimators=50 \\
AdaBoost & learning\_rate=1.0, n\_estimators=100 & learning\_rate=1.0, n\_estimators=100 \\
Gradient Boosting & learning\_rate=1.0, max\_depth=2, n\_estimators=100 & learning\_rate=1.0, max\_depth=2, n\_estimators=100 \\
XGBoost & learning\_rate=0.1, max\_depth=3, n\_estimators=100 & learning\_rate=1.0, max\_depth=3, n\_estimators=50 \\
\bottomrule
\end{longtable}

\subsection*{Table 5: 5-Fold Cross-Validation Results -- No-PCA}
\begin{longtable}{lcccccc}
\toprule
Model & Fold 1 & Fold 2 & Fold 3 & Fold 4 & Fold 5 & Avg (No-PCA)\\
\midrule
SVM & 0.9560 & 0.9670 & 0.9560 & 0.9780 & 0.9890 & 0.9692 \\
Naive Bayes & 0.9121 & 0.9670 & 0.8901 & 0.9560 & 0.9451 & 0.9341 \\
KNN & 0.9670 & 0.9670 & 0.9451 & 0.9560 & 0.9780 & 0.9626 \\
Logistic Regression & 0.9780 & 0.9780 & 0.9890 & 0.9780 & 0.9890 & 0.9824 \\
Decision Tree & 0.9451 & 0.9231 & 0.9121 & 0.9231 & 0.8901 & 0.9187 \\
Random Forest & 0.9670 & 0.9560 & 0.9341 & 0.9670 & 0.9890 & 0.9626 \\
AdaBoost & 0.9890 & 1.0000 & 0.9451 & 0.9780 & 0.9890 & 0.9802 \\
Gradient Boosting & 0.9670 & 0.9670 & 0.9560 & 0.9451 & 0.9780 & 0.9626 \\
XGBoost & 0.9890 & 0.9780 & 0.9451 & 0.9560 & 0.9890 & 0.9714 \\
Stacking & 0.9670 & 0.9670 & 0.9341 & 0.9670 & 0.9890 & 0.9648 \\
\bottomrule
\end{longtable}

\subsection*{Table 5 (cont.): 5-Fold Cross-Validation Results -- With-PCA}
\begin{longtable}{lcccccc}
\toprule
Model & Fold 1 & Fold 2 & Fold 3 & Fold 4 & Fold 5 & Avg (With-PCA)\\
\midrule
SVM & 0.9560 & 0.9780 & 0.9560 & 0.9670 & 0.9890 & 0.9692 \\
Naive Bayes & 0.8791 & 0.9341 & 0.9121 & 0.9341 & 0.9231 & 0.9165 \\
KNN & 0.9670 & 0.9780 & 0.9231 & 0.9451 & 0.9890 & 0.9604 \\
Logistic Regression & 0.9780 & 0.9780 & 0.9890 & 0.9780 & 0.9890 & 0.9824 \\
Decision Tree & 0.9121 & 0.8901 & 0.8901 & 0.9670 & 0.9560 & 0.9231 \\
Random Forest & 0.9451 & 0.9670 & 0.9231 & 0.9780 & 0.9670 & 0.9560 \\
AdaBoost & 0.9231 & 0.9780 & 0.9121 & 0.9670 & 0.9780 & 0.9516 \\
Gradient Boosting & 0.9451 & 0.9780 & 0.9341 & 0.9780 & 0.9780 & 0.9626 \\
XGBoost & 0.9231 & 0.9560 & 0.9560 & 0.9780 & 0.9670 & 0.9560 \\
Stacking & 0.9451 & 0.9780 & 0.9560 & 0.9670 & 0.9780 & 0.9648 \\
\bottomrule
\end{longtable}

\subsection*{Test Set Performance (No-PCA vs With-PCA)}
\begin{longtable}{lcccc}
\toprule
Model & Test Acc (No-PCA) & Test Acc (With-PCA) & Test F1 (No-PCA) & Test F1 (With-PCA)\\
\midrule
SVM & 0.9737 & 0.9737 & 0.9790 & 0.9790 \\
Naive Bayes & 0.9298 & 0.9211 & 0.9444 & 0.9379 \\
KNN & \textbf{0.9825} & 0.9561 & \textbf{0.9863} & 0.9655 \\
Logistic Regression & 0.9737 & 0.9737 & 0.9793 & 0.9793 \\
Decision Tree & 0.9211 & 0.8947 & 0.9362 & 0.9143 \\
Random Forest & 0.9561 & 0.9123 & 0.9655 & 0.9296 \\
AdaBoost & 0.9561 & 0.9386 & 0.9660 & 0.9510 \\
Gradient Boosting & 0.9649 & 0.9386 & 0.9726 & 0.9517 \\
XGBoost & 0.9474 & 0.9561 & 0.9589 & 0.9650 \\
Stacking & 0.9737 & 0.9649 & 0.9793 & 0.9722 \\
\bottomrule
\end{longtable}

\section{Result Visualization}

\begin{figure}[H]
\centering
\includegraphics[width=0.9\linewidth]{accuracy_comparison.png}
\caption{Average 5-Fold CV Accuracy per model: No-PCA vs With-PCA}
\end{figure}

\textbf{Inference:} SVM, Logistic Regression, and Gradient Boosting show virtually
no change in CV accuracy after PCA (differences under 0.1 percentage points),
while Naive Bayes, Random Forest, AdaBoost, and XGBoost all show a visible drop
under PCA. Decision Tree is the only model whose CV accuracy slightly
\emph{improved} with PCA (91.87\% $\rightarrow$ 92.31\%), though its test accuracy
still dropped -- a sign that improvement was fold-specific rather than a genuine
generalization gain.

\begin{figure}[H]
\centering
\includegraphics[width=0.75\linewidth]{roc_curves.png}
\caption{ROC Curves -- SVM and Stacking, No-PCA vs With-PCA}
\end{figure}

\textbf{Inference:} SVM's ROC curve is nearly identical between No-PCA and
With-PCA, consistent with its unchanged accuracy above. Stacking's With-PCA curve
sits marginally below its No-PCA curve, matching the small test-accuracy drop
(97.37\% $\rightarrow$ 96.49\%) seen in the results table -- PCA cost Stacking a
small amount of the diversity its three heterogeneous base learners relied on.

\begin{figure}[H]
\centering
\includegraphics[width=0.45\linewidth]{cm_svm_NoPCA.png}
\includegraphics[width=0.45\linewidth]{cm_svm_WithPCA.png}
\caption{Confusion Matrices -- SVM, No-PCA (left) vs With-PCA (right)}
\end{figure}

\textbf{Inference:} SVM's confusion matrices are effectively identical under both
settings -- the same number of misclassifications in each case -- confirming that
this particular model's decision boundary transfers cleanly into the 10-component
PCA space without losing separative power.

\begin{figure}[H]
\centering
\includegraphics[width=0.45\linewidth]{cm_stacking_NoPCA.png}
\includegraphics[width=0.45\linewidth]{cm_stacking_WithPCA.png}
\caption{Confusion Matrices -- Stacking, No-PCA (left) vs With-PCA (right)}
\end{figure}

\textbf{Inference:} Stacking shows one additional misclassification under
With-PCA compared to No-PCA, consistent with its small test-accuracy and F1 drop
-- a modest but real cost of dimensionality reduction for this particular
ensemble.

\section{Discussion}

\subsection*{Which models improved most with PCA? Which did not? Why?}
No model showed a meaningful \emph{improvement} in test accuracy under PCA in
this experiment; the closest cases were SVM and Logistic Regression, which were
essentially unaffected (identical test accuracy before and after). Tree-based and
boosting models (Random Forest, AdaBoost, Gradient Boosting, Decision Tree) showed
the clearest \emph{drops} under PCA -- these models naturally perform their own
implicit feature selection by choosing informative axis-aligned splits, so
compressing the original, individually-interpretable features into 10 rotated PCA
components removes exactly the kind of structure they exploit best.

\subsection*{Did PCA reduce variance across folds (more stable results)?}
Comparing the fold-to-fold spread in Table 5, PCA did not clearly reduce
variance across folds for most models -- e.g. AdaBoost's fold scores ranged from
0.9121 to 0.9780 with PCA (a wider spread) versus 0.9451 to 1.0000 without. Only
Decision Tree showed a slightly tighter, more stable fold range under PCA. Overall,
stability was comparable between the two settings on this dataset, rather than
PCA producing a clear stabilizing effect.

\subsection*{For high-dimensional data, was PCA beneficial in reducing overfitting?}
Since the original 30-feature space is only moderately high-dimensional relative
to 569 samples, most models were not overfitting heavily to begin with, so PCA had
limited room to help. If anything, PCA slightly \emph{increased} the train-test
gap for models like Random Forest and AdaBoost (test accuracy dropped more than
CV accuracy did), suggesting the compressed feature space made it marginally
harder for these ensembles to generalize rather than easier.

\subsection*{How did linear models (Logistic Regression, SVM) behave compared to ensemble models with PCA?}
Both linear-style models (Logistic Regression and SVM with an RBF kernel) were
essentially unaffected by PCA (test accuracy identical, F1 identical). Ensemble
models were more sensitive: Random Forest, AdaBoost, and Gradient Boosting all
lost 1--3 percentage points of test accuracy under PCA, while XGBoost was a
mild exception, gaining slightly (94.74\% $\rightarrow$ 95.61\%). This suggests
linear decision boundaries in a rotated PCA space are just as easy to draw as in
the original space, while tree-based ensembles lose some of the axis-aligned
structure they depend on.

\subsection*{Did stacking show robustness to dimensionality reduction compared to single models?}
Stacking's test accuracy dropped only slightly under PCA (97.37\% $\rightarrow$
96.49\%, a 0.88-point drop) -- smaller than Random Forest's 4.38-point drop or
AdaBoost's 1.75-point drop, though larger than SVM/Logistic Regression's 0-point
change. This places Stacking in a middle ground: more robust to PCA than most
individual tree-based ensembles (since its SVM and Naive Bayes base learners are
relatively PCA-tolerant), but not as fully unaffected as the purely linear
models.

\section{Conclusion}
On this dataset, PCA reduction to 10 components (95\% variance) neither
consistently helped nor catastrophically hurt performance, but its effect was
clearly \textbf{model-dependent}: linear and margin-based models (Logistic
Regression, SVM) were essentially unaffected, while tree-based and boosting
ensembles (Random Forest, AdaBoost, Gradient Boosting, Decision Tree) lost
meaningful test accuracy. \textbf{KNN achieved the best No-PCA test performance}
(98.25\% accuracy) but also showed one of the larger PCA-related drops (95.61\%),
while \textbf{SVM and Logistic Regression were the most PCA-robust} of the
individually-tuned models. Given the strong feature correlations already observed
in EDA, PCA's main practical benefit here is dimensionality reduction itself (30
$\rightarrow$ 10 features) for faster training and simpler downstream pipelines,
rather than an accuracy improvement -- and that trade-off is worth making mainly
for the models (SVM, Logistic Regression, XGBoost) that tolerate it well, and
worth avoiding for models (Random Forest, AdaBoost, Decision Tree) that rely on
the original feature structure.

\section{References}
\begin{enumerate}
\item Scikit-learn: Decomposing signals in components (PCA), \url{https://scikit-learn.org/stable/modules/decomposition.html}
\item Scikit-learn: Ensemble Methods, \url{https://scikit-learn.org/stable/modules/ensemble.html}
\item XGBoost Documentation, \url{https://xgboost.readthedocs.io}
\item UCI Machine Learning Repository: Breast Cancer Wisconsin (Diagnostic) Dataset
\end{enumerate}

\end{document}
